\section{Mathematical Description of the Padding Oracle Attack}

This section develops the algebraic structure behind CBC padding oracle attacks in detail. The central point is that although the attacker does not know the key and does not observe the plaintext directly, the attacker can still manipulate ciphertext blocks so that the decryption routine evaluates carefully chosen predicates on the resulting plaintext. The padding-validity response then leaks enough information to recover the plaintext block by block.

\subsection{CBC decryption and the intermediate value}

Let the block length be $b$ bytes, and let the ciphertext be written as
\[
C = C_0 \| C_1 \| \cdots \| C_n,
\]
where $C_0$ is the initialization vector and $C_1,\dots,C_n$ are the ciphertext blocks. For $i \ge 1$, CBC decryption computes
\[
P_i = D_k(C_i) \oplus C_{i-1}.
\]
It is convenient to define the \emph{intermediate value}
\[
I_i = D_k(C_i),
\]
so that
\[
P_i = I_i \oplus C_{i-1}.
\]
The attacker does not know $I_i$, but $C_{i-1}$ is fully visible and can be replaced by an attacker-chosen block $C'_{i-1}$. If the target block $C_i$ is left fixed, then decryption of the modified pair $(C'_{i-1}, C_i)$ yields
\[
P'_i = D_k(C_i) \oplus C'_{i-1} = I_i \oplus C'_{i-1}.
\]
Thus the attacker can control the output block $P'_i$ up to the unknown mask $I_i$.

\subsection{Padding validity as a predicate}

Assume the implementation uses PKCS\#7 padding on a $b$-byte block. A string $x \in \{0,1\}^{8b}$ has valid padding of length $r$ if its final $r$ bytes all equal the byte value $r$, where $1 \le r \le b$. For example, valid suffixes include
\[
01, \qquad 02\,02, \qquad 03\,03\,03, \quad \dots
\]
and so on up to $b$ repeated bytes.

Define the padding oracle
\[
\mathcal{O}(X) =
\begin{cases}
1 & \text{if decryption of } X \text{ yields valid PKCS\#7 padding},\\
0 & \text{otherwise.}
\end{cases}
\]
The attacker is given oracle access to $\mathcal{O}$ and a target ciphertext $C_0\|C_1\|\cdots\|C_n$. The objective is to recover the plaintext without the key.

The crucial observation is that valid padding is a structured property of the final bytes of the decrypted block. Since the attacker can vary $C'_{i-1}$ and observe whether $P'_i = I_i \oplus C'_{i-1}$ lies in the set of correctly padded strings, the oracle effectively leaks whether certain linear equations on the bytes of $I_i$ hold.

\subsection{Recovering the last byte}

Consider a fixed target block $C_i$ and let the attacker vary only the final byte of $C'_{i-1}$. Write bytes with index $b$ for the last position. Then
\[
P'_i[b] = I_i[b] \oplus C'_{i-1}[b].
\]
The attacker wants to force the last byte of $P'_i$ to equal $01$, which is valid padding of length one. Therefore the attacker searches for a guess $g \in \{0,\dots,255\}$ such that
\[
I_i[b] \oplus g = 01.
\]
Whenever the oracle returns valid padding, the attacker obtains
\[
I_i[b] = g \oplus 01.
\]
Since the original plaintext byte satisfies
\[
P_i[b] = I_i[b] \oplus C_{i-1}[b],
\]
it follows that
\[
P_i[b] = (g \oplus 01) \oplus C_{i-1}[b].
\]
This recovers the final plaintext byte.

\subsection{Recovering the second-to-last byte}

After recovering $I_i[b]$, the attacker moves to position $b-1$. Now the goal is to force the final two bytes of $P'_i$ to equal $02\,02$. The attacker first sets the last byte of $C'_{i-1}$ so that
\[
P'_i[b] = I_i[b] \oplus C'_{i-1}[b] = 02,
\]
which is achieved by choosing
\[
C'_{i-1}[b] = I_i[b] \oplus 02.
\]
Then the attacker varies the byte at position $b-1$ until the oracle accepts. For the successful guess $g$,
\[
I_i[b-1] \oplus g = 02,
\]
so
\[
I_i[b-1] = g \oplus 02.
\]
Hence
\[
P_i[b-1] = I_i[b-1] \oplus C_{i-1}[b-1] = (g \oplus 02) \oplus C_{i-1}[b-1].
\]

\subsection{General byte-recovery step}

Suppose the attacker has already recovered the final $r-1$ bytes of $I_i$, namely
\[
I_i[b-r+2], I_i[b-r+3], \dots, I_i[b].
\]
To recover the next byte $I_i[b-r+1]$, the attacker wants the final $r$ bytes of $P'_i$ to all equal $r$. For each known suffix position $j \in \{b-r+2,\dots,b\}$, the attacker sets
\[
C'_{i-1}[j] = I_i[j] \oplus r,
\]
which guarantees
\[
P'_i[j] = I_i[j] \oplus C'_{i-1}[j] = r.
\]
Then the attacker varies the unknown position $b-r+1$ until the oracle returns valid padding. If the successful value is $g$, then
\[
I_i[b-r+1] \oplus g = r,
\]
and therefore
\[
I_i[b-r+1] = g \oplus r.
\]
The original plaintext byte is then
\[
P_i[b-r+1] = I_i[b-r+1] \oplus C_{i-1}[b-r+1].
\]
Repeating this process for $r=1,2,\dots,b$ recovers the entire block $P_i$.

\subsection{Why the attack works}

The attack succeeds because CBC decryption is linearly malleable with respect to the previous ciphertext block. The attacker does not need to invert $D_k$ or distinguish the block cipher from random. Instead, the attacker uses the oracle to test whether carefully crafted values of $I_i \oplus C'_{i-1}$ land inside the valid-padding set. Since PKCS\#7 padding provides a nested family of highly structured suffix constraints, each successful oracle response reveals one more byte of the intermediate value $I_i$.

Equivalently, the oracle converts the hidden value $I_i$ into something searchable one byte at a time. The attacker cannot ask directly for $I_i[j]$, but can ask whether
\[
I_i[j] \oplus g = r
\]
for chosen values of $g$ and $r$. This is enough to solve for the byte.

\subsection{Complexity estimate}

In the worst case, recovering one byte requires trying all $256$ possibilities. For a block of $b$ bytes, this yields an upper bound of roughly
\[
256b
\]
oracle queries per block, ignoring small overheads due to false-positive handling. For AES with $b=16$, this is about
\[
256 \cdot 16 = 4096
\]
queries per block in the simplest analysis. In practice the expected number is smaller because the correct guess is found on average halfway through the search. The expected cost is therefore closer to
\[
128b
\]
queries per block, plus a small constant-factor overhead.

If the message contains $m$ ciphertext blocks, then the total attack complexity is linear in $m$ and linear in the block size. This makes the attack entirely practical whenever an adversary can submit a moderate number of oracle queries.

\subsection{Worked symbolic example}

To make the derivation concrete, consider a two-block ciphertext $C_0\|C_1$ and suppose the attacker wants to recover $P_1$. The attacker fixes $C_1$ and submits modified IVs $C'_0$. Since
\[
P'_1 = D_k(C_1) \oplus C'_0 = I_1 \oplus C'_0,
\]
valid padding of length one occurs exactly when
\[
I_1[16] \oplus C'_0[16] = 01.
\]
If the successful final-byte guess is $g$, then
\[
I_1[16] = g \oplus 01,
\qquad
P_1[16] = I_1[16] \oplus C_0[16].
\]
Next the attacker sets
\[
C'_0[16] = I_1[16] \oplus 02
\]
and varies $C'_0[15]$ until the oracle accepts. If the successful guess is $h$, then
\[
I_1[15] = h \oplus 02,
\qquad
P_1[15] = I_1[15] \oplus C_0[15].
\]
Continuing in this way recovers all of $P_1$.

\subsection{Security interpretation}

This derivation makes clear that the attack does not break AES as a primitive. The problem is that CBC with padding but without integrity protection permits controlled modification of the decrypted plaintext, while the implementation leaks whether the modified plaintext satisfies a recognizable syntactic condition. In other words, confidentiality is lost not because the block cipher is weak, but because the scheme allows malleability and then exposes information about the result of decryption.

This is precisely why authenticated encryption fixes the problem. Once ciphertext integrity is checked before releasing any information about the decrypted plaintext, the adversary no longer obtains the yes/no feedback needed to solve for the intermediate values. The padding oracle disappears because manipulated ciphertexts are rejected at the authentication layer.

\section{Coding Demo and Experimental Component}

To complement the mathematical analysis, this project includes a local Python demonstration of the attack. The code constructs a toy AES-CBC system with PKCS\#7 padding and exposes a deliberately insecure oracle that reveals only whether padding is valid. An attack script then recovers the plaintext using oracle access alone.

The implementation follows the mathematical derivation above. For each target block $C_i$, the code keeps $C_i$ fixed and replaces the previous block by an attacker-controlled value $C'_{i-1}$. It first forces a valid padding suffix of length one, then of length two, and so on until the entire block is recovered. The code also records the total number of oracle queries, which provides an empirical comparison with the theoretical complexity estimate.

To make the security lesson explicit, the code also includes a comparison with AES-GCM. In that setting, modifying the ciphertext causes authentication failure before any padding-style feedback is exposed. This side-by-side comparison highlights the central lesson of the report: CBC padding oracle attacks are not merely about one unfortunate implementation detail, but about the need to combine privacy with integrity.

The current implementation is complete enough for the interim report because it already demonstrates the vulnerable design, the recovery algorithm, and the authenticated-encryption contrast. Remaining work for the final report includes polishing the presentation of experimental results, possibly adding a small table of query counts over several runs, and integrating screenshots or terminal output into the final writeup.
